% !TeX root = ../../../main.tex
\section{问题重述}
\subsection{研究背景}
这里填写题目所涉及的实际背景和研究目标。正文采用宋体风格中文与 TeX Gyre Termes 西文搭配，段首缩进两个汉字，行距与图表标题参考 ElegantPaper。本文中的数值、文献和示例图仅用于展示排版效果。

\subsection{待解决的问题}
\begin{enumerate}
  \item 问题一：填写题目要求完成的任务、输入数据与预期输出。
  \item 问题二：说明与上一问的联系及新增约束。
  \item 问题三：填写结果检验、优化或决策分析要求。
\end{enumerate}

\section{问题分析}
\subsection{整体思路}
在此说明数据处理、建模、求解与检验的关系，并解释各问题之间如何衔接。示例书籍引用应注明具体页码\cite[第23--25页]{samplebook}；参考文献编号由首次引用顺序自动生成。

\subsection{各问题的分析}
分别解释每一问需要解决的核心矛盾、选用模型的理由及其适用条件。不要将尚未完成的实验或模型效果写成已经验证的结论。
